\documentclass[border=6mm]{standalone}
% ==============================================================================
% SETUP.TEX - CONFIGURACIÓN GLOBAL DEL PROYECTO
% ==============================================================================
% Este archivo centraliza todos los paquetes y estilos.
% Debe incluirse en el preámbulo del libro principal y de los archivos standalone.
% \PassOptionsToPackage{gray}{xcolor}
% --- 1. CODIFICACIÓN E IDIOMA ---
% fix-cm habilita las fuentes Computer Modern en tamaños arbitrarios (por debajo
% de 5 pt, entre otros). Debe cargarse antes que fontenc.
\usepackage{fix-cm}
\usepackage[utf8]{inputenc}
\usepackage[T1]{fontenc}
\usepackage[spanish, es-tabla]{babel}  
\usepackage{csquotes} % Recomendado para babel y citas

\usepackage{import}
% mode=image: Usa el PDF pre-compilado, nunca intenta recompilar.
% Debes compilar cada figura manualmente antes de compilar el libro.
\usepackage[mode=image, subpreambles=false]{standalone}

% --- 2. MATEMÁTICAS Y CIENCIAS ---
\usepackage{amsmath, amssymb, amsfonts, mathtools}
\usepackage{enumitem}

% LaTeX trae \DeclareMathSizes{5}{5}{5}{5}, así que a 5 pt (\tiny) los subíndices
% salen del mismo tamaño que la base: en las etiquetas de figuras el M de
% $\omega_M$ queda desproporcionado. Se restablece la proporción habitual.
\DeclareMathSizes{5}{5}{3.7}{3}

% --- 3. DISEÑO Y GEOMETRÍA --- 
% Unificamos los márgenes para todo el documento
\usepackage[left=2.5cm,right=2.5cm,top=3cm,bottom=3cm]{geometry}
\makeatletter
\ifstandalone % Evita que geometry dañe el recorte de las figuras bajándolas 3cm
  \@ifclasswith{standalone}{crop=false}{
    % Es un capítulo/sección: Dejamos que geometry actúe.
  }{
    % Es una figura: Neutralizamos geometry para que no las recorte mal.
    \geometry{pass}
  }
\fi
\makeatother
\usepackage{parskip} % Espaciado entre párrafos sin sangría
\usepackage{float}   % Para usar [H] en figura

% Ajustes de flotantes para evitar espacios verticales exagerados
\renewcommand{\topfraction}{0.95}      % Permite que las figuras ocupen hasta el 95% superior
\renewcommand{\bottomfraction}{0.95}   % Permite que las figuras ocupen hasta el 95% inferior
\renewcommand{\textfraction}{0.05}     % Permite hojas con tan solo un 5% de texto
\renewcommand{\floatpagefraction}{0.8} % Requiere que una página de solo figuras esté 80% llena
\setcounter{topnumber}{4}
\setcounter{bottomnumber}{4}
\setcounter{totalnumber}{10}

% --- 4. GRÁFICOS Y TABLAS ---
\usepackage{graphicx}
\usepackage{booktabs} % Tablas profesionales
\usepackage{caption}  % Control de captions y \captionof
\usepackage{tikz}
\usetikzlibrary{arrows.meta, calc, angles, quotes, babel, backgrounds}
\usepackage{pgfplots}
\usepgfplotslibrary{groupplots}
\usepgfplotslibrary{external}
\usepgfplotslibrary{patchplots}
\pgfplotsset{compat=1.18}
\usepackage[american, straightvoltages]{circuitikz}

% --- 5. COLORES Y CAJAS ---

\usepackage[table]{xcolor}
\usepackage{tcolorbox}

% Definición de colores corporativos/estilo
\definecolor{utnblue}{RGB}{0, 51, 102}
\definecolor{codegreen}{rgb}{0,0.6,0}
\definecolor{codegray}{rgb}{0.5,0.5,0.5}
\definecolor{codepurple}{rgb}{0.58,0,0.82}
\definecolor{backcolour}{rgb}{0.96,0.96,0.96}

% --- 6. ENTORNOS PERSONALIZADOS (CAJAS) ---

% Caja "Resultado" (Azul Corporativo) — Definiciones, fórmulas clave, resultados importantes.
% Uso: \begin{resultado}[Fórmula de Euler] ... \end{resultado}
\newtcolorbox{resultado}[1][]{
    colback=utnblue!5!white,
    colframe=utnblue,
    title=\textbf{#1},
    fonttitle=\bfseries,
    sharp corners,
    boxrule=0.5mm
}

% Caja "Nota" (Gris) — Advertencias, aplicaciones, contexto secundario.
% Uso: \begin{nota}[Cuidado con la calculadora] ... \end{nota}
% Uso: \begin{nota} ... \end{nota}  → título por defecto "Nota"
\newtcolorbox{nota}[1][Nota]{
    colback=codegray!5!white,
    colframe=codegray,
    title=\textbf{#1},
    fonttitle=\bfseries,
    sharp corners,
    boxrule=0.5mm
}

% Caja inline para resaltar un valor y enlazarlo con su otra aparición en una tabla
% o en una cuenta: mismo color, mismo valor.
% Uso: \marcavalor{$4$}  o  \marcavalor[codegreen]{$4$}
\newtcbox{\marcavalor}[1][utnblue]{
    on line,
    boxsep=1pt, left=2pt, right=2pt, top=1pt, bottom=1pt,
    colback=#1!10!white,
    colframe=#1,
    boxrule=0.4pt,
    arc=2pt
}

% --- 7. CÓDIGO FUENTE (LISTINGS) ---
\usepackage{listings}

\lstdefinestyle{miestilo}{
    backgroundcolor=\color{backcolour},   
    commentstyle=\color{codegreen},
    keywordstyle=\color{magenta},
    numberstyle=\tiny\color{codegray},
    stringstyle=\color{codepurple},
    basicstyle=\ttfamily\footnotesize,
    breakatwhitespace=false,         
    breaklines=true,                 
    captionpos=b,                    
    keepspaces=true,                 
    numbers=left,                    
    numbersep=5pt,                  
    showspaces=false,                
    showstringspaces=false,
    showtabs=false,                  
    tabsize=2,
    frame=single,                    
    rulecolor=\color{codegray}
}

\lstset{style=miestilo}
% Soporte para tildes y ñ en bloques de código
\lstset{literate={ñ}{{\~n}}1 {á}{{\'a}}1 {é}{{\'e}}1 {í}{{\'i}}1 {ó}{{\'o}}1 {ú}{{\'u}}1}

% Operadores matemáticos personalizados

% Vector en sentido discreto (lista finita de coordenadas): negrita + flecha.
% Uso: \vect{v}, \vect{e}_k, \vect{0}.
\newcommand{\vect}[1]{\vec{\mathbf{#1}}}

% Fecha de versión y comando para capítulos standalone
\providecommand{\verdate}{\the\year.\ifnum\month<10 0\fi\the\month.\ifnum\day<10 0\fi\the\day}
\newcommand{\chapterversion}{%
    \vspace{-1.5cm}\begin{flushright}\small\textit{\color{codegray}Versión~\verdate}\end{flushright}\vspace{0.5cm}%
}

% --- 8. HIPERVÍNCULOS (DEBE SER EL ÚLTIMO PAQUETE) ---
\usepackage[colorlinks=true, linkcolor=utnblue, urlcolor=utnblue, citecolor=utnblue]{hyperref}

% --- 9. REFERENCIAS CRUZADAS ENTRE CAPÍTULOS STANDALONE ---
% Cuando se compila un capítulo standalone, xr-hyper lee los .aux de los
% otros capítulos (si existen) para resolver \ref{} a labels externos.
% En la compilación del libro completo este bloque no se activa.
\ifstandalone
  \usepackage{xr-hyper}
  \makeatletter
  \newcommand{\xrchap}[1]{%
    \edef\xrc@self{\detokenize{Capitulo#1}}%
    \edef\xrc@job{\jobname}%
    \ifx\xrc@self\xrc@job\else
      \IfFileExists{../#1capitulo/Capitulo#1.aux}%
        {\externaldocument{../#1capitulo/Capitulo#1}}{}%
    \fi
  }
  \makeatother
  \xrchap{01}\xrchap{02}\xrchap{03}\xrchap{04}
  \xrchap{05}\xrchap{06}\xrchap{07}\xrchap{08}
  \xrchap{09}\xrchap{10}\xrchap{11}\xrchap{12}
  \xrchap{13}
\fi
\usetikzlibrary{shadows.blur, shapes.geometric, calc}

% Realización directa con derivadores: caso general.
% EDLCC:  sum_{k=0}^{N} a_k y^(k) = sum_{k=0}^{M} b_k x^(k)
% Despejada: y = (1/a_0)[ sum b_k x^(k) - sum_{k>=1} a_k y^(k) ]
%
% Estructura:
%   - Cadena de M derivadores en la columna de la entrada (izquierda).
%   - Cadena de N derivadores en la columna de la salida (derecha).
%   - Escaladores b_k/a_0 alineados en una columna entre la entrada y los
%     sumadores, con el triángulo apuntando a la derecha.
%   - Escaladores -a_k/a_0 alineados en una columna entre los sumadores y
%     la salida, con el triángulo apuntando a la izquierda.
%   - Cadena de sumadores en la columna central: cada uno junta la
%     contribución b_k x^(k) (oeste), la contribución -a_k y^(k) (este) y
%     la suma parcial proveniente de la fila inferior (sur). La salida (norte)
%     alimenta el sumador de la fila superior. El sumador superior produce y(t).
%   - Entre el segundo y el último elemento de cada cadena se dibuja una
%     flecha discontinua que sintetiza los términos intermedios.

\begin{document}
\begin{tikzpicture}[>=Latex, font=\small,
    gain_r/.style={
        regular polygon, regular polygon sides=3, shape border rotate=30,
        draw=utnblue!60, line width=0.5pt, fill=utnblue!8,
        minimum size=1.35cm, inner sep=0pt, font=\scriptsize,
        blur shadow={shadow blur steps=4, shadow blur extra rounding=0pt,
                     shadow xshift=0.4pt, shadow yshift=-0.6pt, shadow opacity=25}
    },
    gain_l/.style={
        regular polygon, regular polygon sides=3, shape border rotate=210,
        draw=utnblue!60, line width=0.5pt, fill=utnblue!8,
        minimum size=1.35cm, inner sep=0pt, font=\scriptsize,
        blur shadow={shadow blur steps=4, shadow blur extra rounding=0pt,
                     shadow xshift=0.4pt, shadow yshift=-0.6pt, shadow opacity=25}
    },
    op_block/.style={
        draw=utnblue!60, line width=0.5pt, fill=utnblue!8,
        minimum width=1.2cm, minimum height=0.9cm, font=\normalsize,
        blur shadow={shadow blur steps=4, shadow blur extra rounding=0pt,
                     shadow xshift=0.4pt, shadow yshift=-0.6pt, shadow opacity=25}
    },
    sum/.style={
        draw=utnblue!60, line width=0.5pt, circle,
        minimum size=0.9cm, inner sep=0pt, fill=white,
        blur shadow={shadow blur steps=4, shadow blur extra rounding=0pt,
                     shadow xshift=0.3pt, shadow yshift=-0.4pt, shadow opacity=20}
    },
    cross_line/.style={utnblue!40, line width=0.4pt},
    signal/.style={->, thick, utnblue},
    wire/.style={-, thick, utnblue},
    dashsignal/.style={->, thick, utnblue, dashed},
    dashwire/.style={-, thick, utnblue, dashed},
    tlabel/.style={font=\footnotesize, inner sep=2pt},
]

% --- Coordenadas ---
\def\xinL{-7.2}
\def\xin{-5.6}
\def\xb{-2.8}
\def\xs{0}
\def\xa{2.8}
\def\xout{5.6}
\def\xoutR{7.2}

\def\yzero{0}
\def\ydone{-1.1}
\def\ytone{-2.2}
\def\ydtwo{-3.3}
\def\yttwo{-4.4}
\def\ydlast{-6.4}
\def\ytlast{-7.5}

% =========================================================================
% Macros locales para dibujar las cruces y signos del sumador.
% \sumcrossWS{node}   -> cruz + signos en O y S  (sumador superior: salida E)
% \sumcrossWES{node}  -> cruz + signos en O, E y S (sumador intermedio: salida N)
% \sumcrossWE{node}   -> cruz + signos en O y E  (sumador inferior: salida N)
% =========================================================================
\newcommand{\sumcrossbase}[1]{%
    \draw[cross_line] ([xshift=-0.318cm,yshift=-0.318cm]#1.center) -- ([xshift=0.318cm,yshift=0.318cm]#1.center);
    \draw[cross_line] ([xshift=-0.318cm,yshift=0.318cm]#1.center) -- ([xshift=0.318cm,yshift=-0.318cm]#1.center);
}
\newcommand{\sumcrossWS}[1]{%
    \sumcrossbase{#1}%
    \node[font=\scriptsize] at ([xshift=-0.27cm]#1.center) {$+$};%
    \node[font=\scriptsize] at ([yshift=-0.27cm]#1.center) {$+$};%
}
\newcommand{\sumcrossWES}[1]{%
    \sumcrossbase{#1}%
    \node[font=\scriptsize] at ([xshift=-0.27cm]#1.center) {$+$};%
    \node[font=\scriptsize] at ([xshift=0.27cm]#1.center)  {$+$};%
    \node[font=\scriptsize] at ([yshift=-0.27cm]#1.center) {$+$};%
}
\newcommand{\sumcrossWE}[1]{%
    \sumcrossbase{#1}%
    \node[font=\scriptsize] at ([xshift=-0.27cm]#1.center) {$+$};%
    \node[font=\scriptsize] at ([xshift=0.27cm]#1.center)  {$+$};%
}

% =========================================================================
% Cadena de sumadores en la columna central
% =========================================================================
\node[sum] (S0)    at (\xs, \yzero)  {};   \sumcrossWS{S0}
\node[sum] (S1)    at (\xs, \ytone)  {};   \sumcrossWES{S1}
\node[sum] (S2)    at (\xs, \yttwo)  {};   \sumcrossWES{S2}
\node[sum] (Slast) at (\xs, \ytlast) {};   \sumcrossWE{Slast}

% =========================================================================
% Espina superior: x(t) -> tap -> [b0/a0] -> S0 -> tap -> y(t)
% =========================================================================
\coordinate (in_left)   at (\xinL, \yzero);
\coordinate (in_tap)    at (\xin,  \yzero);
\coordinate (out_tap)   at (\xout, \yzero);
\coordinate (out_right) at (\xoutR,\yzero);

\node[gain_r] (B0) at (\xb, \yzero) {$\dfrac{b_0}{a_0}$};

\draw[wire]   (in_left) -- (in_tap) node[above, midway] {$x(t)$};
\fill[utnblue] (in_tap) circle (2.2pt);
\draw[signal] (in_tap) -- (B0.west); 
\draw[signal] (B0.east) -- (S0.west);
\draw[wire]   (S0.east) -- (out_tap);
\fill[utnblue] (out_tap) circle (2.2pt);
\draw[signal] (out_tap) -- (out_right) node[above, midway] {$y(t)$};

% =========================================================================
% Cadena de derivadores en la entrada (columna izquierda)
% =========================================================================
\node[op_block] (Din1) at (\xin, \ydone) {$\dfrac{d}{dt}$};
\coordinate (tin1) at (\xin, \ytone);
\node[op_block] (Din2) at (\xin, \ydtwo) {$\dfrac{d}{dt}$};
\coordinate (tin2) at (\xin, \yttwo);
\node[op_block] (DinM) at (\xin, \ydlast) {$\dfrac{d}{dt}$};
\coordinate (tinM) at (\xin, \ytlast);

\draw[signal] (in_tap) -- (Din1.north);
\draw[wire]   (Din1.south) -- (tin1)
              node[tlabel, midway, anchor=east, xshift=-2pt] {$\dot x$};
\fill[utnblue] (tin1) circle (2.2pt);
\draw[signal] (tin1) -- (Din2.north);
\draw[wire]   (Din2.south) -- (tin2)
              node[tlabel, midway, anchor=east, xshift=-2pt] {$\ddot x$};
\fill[utnblue] (tin2) circle (2.2pt);
\draw[dashsignal] (tin2) -- (DinM.north);
\draw[wire]   (DinM.south) -- (tinM)
              node[tlabel, midway, anchor=east, xshift=-2pt] {$x^{(M)}$};
% tinM es esquina (DinM -> BM): no necesita nodo

% =========================================================================
% Cadena de derivadores en la salida (columna derecha)
% =========================================================================
\node[op_block] (Dout1) at (\xout, \ydone) {$\dfrac{d}{dt}$};
\coordinate (tout1) at (\xout, \ytone);
\node[op_block] (Dout2) at (\xout, \ydtwo) {$\dfrac{d}{dt}$};
\coordinate (tout2) at (\xout, \yttwo);
\node[op_block] (DoutN) at (\xout, \ydlast) {$\dfrac{d}{dt}$};
\coordinate (toutN) at (\xout, \ytlast);

\draw[signal] (out_tap) -- (Dout1.north);
\draw[wire]   (Dout1.south) -- (tout1)
              node[tlabel, midway, anchor=west, xshift=2pt] {$\dot y$};
\fill[utnblue] (tout1) circle (2.2pt);
\draw[signal] (tout1) -- (Dout2.north);
\draw[wire]   (Dout2.south) -- (tout2)
              node[tlabel, midway, anchor=west, xshift=2pt] {$\ddot y$};
\fill[utnblue] (tout2) circle (2.2pt);
\draw[dashsignal] (tout2) -- (DoutN.north);
\draw[wire]   (DoutN.south) -- (toutN)
              node[tlabel, midway, anchor=west, xshift=2pt] {$y^{(N)}$};
% toutN es esquina (DoutN -> AN): no necesita nodo

% =========================================================================
% Escaladores de la entrada (columna xb), apuntan a la derecha.
% Cada uno entra al sumador correspondiente por el oeste.
% =========================================================================
\node[gain_r] (B1) at (\xb, \ytone)  {$\dfrac{b_1}{a_0}$};
\node[gain_r] (B2) at (\xb, \yttwo)  {$\dfrac{b_2}{a_0}$};
\node[gain_r] (BM) at (\xb, \ytlast) {$\dfrac{b_M}{a_0}$};

\draw[signal] (tin1) -- (B1.west);
\draw[signal] (B1.east) -- (S1.west);
\draw[signal] (tin2) -- (B2.west);
\draw[signal] (B2.east) -- (S2.west);
\draw[signal] (tinM) -- (BM.west);
\draw[signal] (BM.east) -- (Slast.west);

% =========================================================================
% Escaladores de la salida (columna xa), apuntan a la izquierda.
% Cada uno entra al sumador correspondiente por el este.
% =========================================================================
\node[gain_l] (A1) at (\xa, \ytone)  {$-\dfrac{a_1}{a_0}$};
\node[gain_l] (A2) at (\xa, \yttwo)  {$-\dfrac{a_2}{a_0}$};
\node[gain_l] (AN) at (\xa, \ytlast) {$-\dfrac{a_N}{a_0}$};

\draw[signal] (tout1) -- (A1.east);
\draw[signal] (A1.west) -- (S1.east);
\draw[signal] (tout2) -- (A2.east);
\draw[signal] (A2.west) -- (S2.east);
\draw[signal] (toutN) -- (AN.east);
\draw[signal] (AN.west) -- (Slast.east);

% =========================================================================
% Cascada vertical de sumadores: cada Sk.north -> S(k-1).south
% =========================================================================
\draw[signal]     (S1.north) -- (S0.south);
\draw[signal]     (S2.north) -- (S1.south);
\draw[dashsignal] (Slast.north) -- (S2.south);

\end{tikzpicture}
\end{document}
